% ============================================================
% SECTION 05 — FALSIFICATION CRITERION
% Repository: Monette Research — Entropic Gravity & PEIG Framework
% Part of: 01_PHYSICS_CORE/A_ENTROPIC_GRAVITY/master_paper.tex
% Status: Peer-review draft — Phase 2 corrections applied
% Last Updated: March 2026
% ============================================================
%
% PHASE 2 CORRECTIONS APPLIED IN THIS FILE:
%   [C14] Unifies the three slightly different falsification statements
%         appearing across the corpus into one canonical version
%   [C15] Adds the prediction/confirmation threshold (kappa ~ 10^-4)
%         alongside the falsification threshold
%   [C16] Provides the current experimental bounds table with full
%         citations and clear statement of what each bounds
% ============================================================

\section{Falsification Criterion}
\label{sec:falsification}

\subsection{The Centrality of Falsifiability}
\label{subsec:falsifiability_centrality}

A theoretical framework in physics earns the right to be taken seriously
in proportion to how clearly and precisely it can be proven wrong. A
framework that can accommodate any experimental outcome by adjusting
free parameters is not science --- it is numerology. The framework
presented in this paper has one free parameter: the dimensionless
coupling constant $\kappatilde$. This parameter is constrained by the
first-principles estimate $\kappatilde = -(1/4)\alphascreen$ (with
$\alphascreen\in[10^{-4},10^{-2}]$) and bounded from above by existing
experiments.

This section states, with full precision, the experimental conditions
under which this framework would be:
\begin{enumerate}
    \item \textbf{Falsified} --- shown to be incompatible with observation
    at the predicted sensitivity level
    \item \textbf{Confirmed} --- shown to be consistent with a positive
    detection at the predicted coupling level
    \item \textbf{Constrained} --- bounded more tightly than existing
    limits without reaching full confirmation or falsification
\end{enumerate}

\subsection{The Canonical Falsification Criterion}
\label{subsec:canonical_falsification}

\begin{framed}
\noindent\textbf{Canonical Falsification Criterion (v1.0):}

This framework is \textbf{falsified for laboratory-scale relevance}
if all of the following conditions are simultaneously satisfied:

\begin{enumerate}
    \item \textbf{System size}: Macroscopic quantum-coherent ensembles
    with $N \geq 10^6$ entangled qubits (or atoms in equivalent
    entangled states) are used as sources

    \item \textbf{Null result}: No anomalous stress-energy contribution
    is detected beyond the predictions of standard quantum mechanics and
    general relativity with $\kappatilde = 0$

    \item \textbf{Sensitivity}: The measurement sensitivity reaches
    $\Delta p < 10^{-6}$~Pa (equivalently, $\delta|\kappatilde|
    < 10^{-15}$)

    \item \textbf{Statistical power}: At least $n_{\mathrm{runs}}
    \geq 1000$ experimental runs are completed

    \item \textbf{Platform independence}: The null result is reproduced
    across at least two independent experimental platforms from the set
    \{trapped ions, superconducting circuits, neutral atom arrays,
    optomechanical systems\}
\end{enumerate}

\noindent\textbf{Falsification conclusion}: If all five conditions are
satisfied, then $|\kappatilde| < 10^{-15}$, and this framework predicts
no detectable entanglement-gravity coupling for any laboratory-scale
system achievable with foreseeable technology. The framework's relevance
to laboratory-scale gravity engineering is falsified.

\noindent\textbf{What is not falsified}: A null result at the above
sensitivity does not falsify the underlying hypothesis that entanglement
entropy sources curvature at causal horizons (Jacobson's thermodynamic
gravity). It falsifies only the extrapolation of that mechanism to
laboratory scales.
\end{framed}

\subsection{The Confirmation Threshold}
\label{subsec:confirmation}

The framework is \textbf{confirmed at the predicted level} if:

\begin{enumerate}
    \item A differential acceleration $\Delta a(R)$ is detected at
    $\geq 5\sigma$ statistical significance

    \item The measured $\kappatilde$ (extracted via
    Eq.~\eqref{eq:kappa_extraction}) falls within the predicted range:
    \begin{equation}
        \kappatilde \in \left[-2.5\times10^{-3},\;
        -2.5\times10^{-5}\right]
        \label{eq:confirmation_range}
    \end{equation}

    \item The spatial scaling of $\Delta a$ is consistent with the
    predicted $1/R$ dependence from Eq.~\eqref{eq:delta_a}

    \item The result is reproduced independently on at least two
    platforms
\end{enumerate}

\noindent Detection of $|\kappatilde| \sim 10^{-4}$ would constitute
the discovery of a new gravitational coupling. Detection of
$|\kappatilde| \sim 10^{-5}$ to $10^{-6}$ would be a weaker but
still significant confirmation. Any detection at these levels would
constitute one of the most significant results in experimental
gravitational physics since the first direct detection of gravitational
waves~\cite{ligo2016}.

\subsection{The Constraint Regime}
\label{subsec:constraints}

Between full confirmation and full falsification lies the \textbf{constraint
regime}: experimental results that bound $|\kappatilde|$ more tightly
than existing limits without reaching the predicted signal level or the
falsification threshold.

\begin{equation}
    10^{-15} < |\kappatilde|_{\mathrm{bound}} < 10^{-10}
    \quad \Rightarrow \quad \text{constrained, not falsified or confirmed}
    \label{eq:constraint_regime}
\end{equation}

Results in this regime are scientifically valuable: each factor of 10
improvement in the bound progressively constrains the allowed range of
$\alphascreen$ and narrows the class of decoherence models consistent
with the framework.

\subsection{Current Experimental Bounds}
\label{subsec:current_bounds}

The following table summarizes current upper bounds on $|\kappatilde|$
from experiments that are sensitive to anomalous gravitational effects
in quantum systems. All experiments are null results --- no anomalous
coupling has been detected. The bounds are derived by computing the
expected signal from Eq.~\eqref{eq:delta_a} at each experiment's
entanglement density and sensitivity, and requiring the predicted signal
to be below the experimental noise floor.

\begin{table}[H]
    \centering
    \caption{Current experimental upper bounds on $|\kappatilde|$ from
    null results. All values are 95\% confidence level upper bounds.
    The coupling sensitivity column shows the minimum $|\kappatilde|$
    detectable given the system size and measurement sensitivity of each
    experiment.}
    \label{tab:bounds}
    \begin{tabular}{p{4.5cm}p{2.5cm}cc}
        \toprule
        \textbf{Experiment} & \textbf{Reference}
        & \textbf{Bound on $|\kappatilde|$}
        & \textbf{System} \\
        \midrule
        Gravity-mediated entanglement
        & Bose et al.~\cite{bose2023}
        & $< 3\times10^{-9}$
        & Nanodiamond NV centers \\
        Atom interferometry
        & Kasevich et al.~\cite{kasevich2023}
        & $< 1.2\times10^{-10}$
        & $^{87}$Rb BEC \\
        Equivalence principle
        & MICROSCOPE~\cite{microscope2022}
        & $< 8\times10^{-11}$
        & Macroscopic test masses \\
        \midrule
        \textbf{This work (projected)}
        & ---
        & $\delta|\kappatilde| = 3.7\times10^{-13}$
        & $^{87}$Rb GHZ, $N=10^6$ \\
        \textbf{Falsification threshold}
        & ---
        & $|\kappatilde| < 10^{-15}$
        & Multiple platforms \\
        \bottomrule
    \end{tabular}
\end{table}

\noindent\textbf{Reading the table:}
\begin{itemize}
    \item The current best bound ($|\kappatilde| < 8\times10^{-11}$ from
    MICROSCOPE) is derived from equivalence principle tests, not from
    direct entanglement manipulation. It bounds $|\kappatilde|$ under
    the assumption that macroscopic test masses have some residual
    entanglement entropy density --- a conservative and model-dependent
    assumption.

    \item The Kasevich et al. bound ($|\kappatilde| < 1.2\times10^{-10}$)
    is the most directly relevant: it uses atom interferometry on Bose-
    Einstein condensates, the same technique proposed here, but with
    lower entanglement entropy density and shorter coherence times.

    \item The proposed experiment sensitivity ($3.7\times10^{-13}$)
    improves on the best current bound by approximately two orders of
    magnitude --- sufficient to begin probing the lower end of the
    predicted coupling range.

    \item The gap between the current best bound ($10^{-10}$) and the
    predicted realistic coupling ($\sim10^{-5}$ to $10^{-4}$) is
    \emph{five orders of magnitude}. This is large, which means the
    framework's predictions are not in tension with any existing
    experiment. It also means that a definitive test requires a
    measurement $10^5$ times more sensitive than current limits ---
    achievable only with the $N = 10^6$ GHZ entanglement protocol
    proposed here.
\end{itemize}

\subsection{Decision Tree for Experimental Outcomes}
\label{subsec:decision_tree}

The following decision tree maps every possible experimental outcome to
its scientific interpretation:

\begin{enumerate}
    \item \textbf{Detection at $|\kappatilde|\in[10^{-5},10^{-3}]$,
    $1/R$ scaling confirmed, reproduced on 2+ platforms}
    $\Rightarrow$ \textbf{Framework confirmed.} Proceed to engineering
    implications; design second-generation experiments to measure
    $\alphascreen$ as a function of decoherence environment.

    \item \textbf{Detection at $|\kappatilde|\in[10^{-13},10^{-5}]$,
    scaling unclear, not yet reproduced}
    $\Rightarrow$ \textbf{Candidate detection.} Requires replication;
    systematic effects must be ruled out; scaling law must be established.

    \item \textbf{Null result at $\delta|\kappatilde|\sim10^{-12}$ to
    $10^{-10}$}
    $\Rightarrow$ \textbf{New constraint.} Tightens bound on $\kappatilde$;
    requires $\alphascreen < 4\times\delta|\kappatilde|$; not yet
    falsified; proceed to higher-sensitivity platform.

    \item \textbf{Null result at $\delta|\kappatilde| < 10^{-15}$,
    $\geq 1000$ runs, $\geq 2$ platforms, $N\geq 10^6$}
    $\Rightarrow$ \textbf{Framework falsified for laboratory-scale
    relevance.} The ideal coupling $\kappatilde_{\mathrm{ideal}} = -1/4$
    is suppressed by $\alphascreen < 4\times10^{-15}$, rendering any
    laboratory-scale effect undetectable with any foreseeable technology.
    The horizon-scale Jacobson framework remains valid and unaffected.

    \item \textbf{Anomalous signal not consistent with $1/R$ scaling}
    $\Rightarrow$ \textbf{New physics, but not this framework.}
    A different source of anomalous gravity is present; full systematic
    investigation required.
\end{enumerate}

\subsection{Timeline}
\label{subsec:timeline}

Based on current experimental capabilities and the timelines of
comparable atom interferometry experiments:

\begin{itemize}
    \item \textbf{Months 1--6}: Trap characterization, GHZ state
    preparation for $N \leq 10^4$, interferometer calibration, SPAM
    baseline measurement

    \item \textbf{Months 6--12}: Scale to $N = 10^5$; first measurement
    of $\Delta a$ at $R = 10$~cm; establish systematic floor

    \item \textbf{Months 12--18}: Scale to $N = 10^6$ GHZ state (primary
    technical challenge); measure $\Delta a$ at three radii; complete
    1000 runs at each radius

    \item \textbf{Months 18--24}: Statistical analysis, systematic review,
    independent replication on second platform; publication of final result
\end{itemize}

\noindent\textbf{The 24-month timeline is realistic} given current
state-of-the-art in atomic physics. The primary technical bottleneck is
achieving $N = 10^6$ atom GHZ state coherence times $\tau_{\mathrm{coh}}
\geq 1$~s. This is at the frontier of current experimental capability
and may require intermediate milestones at $N = 10^4$ and $N = 10^5$
with correspondingly lower sensitivity.

\subsection{What Either Outcome Means for Physics}
\label{subsec:significance}

\begin{itemize}
    \item \textbf{If confirmed}: The first experimental demonstration that
    quantum entanglement sources macroscopic spacetime curvature beyond
    the predictions of standard GR. Opens a new field of ``entanglement
    engineering for spacetime control.'' Requires extension of the
    Jacobson framework to volume-law entanglement. Immediately motivates
    new experiments on larger coherent systems and different entanglement
    types.

    \item \textbf{If falsified}: Establishes that the thermodynamic
    gravity mechanism is confined to causal horizons and does not extend
    to laboratory-scale quantum systems. Provides a precise quantitative
    boundary for where quantum information physics ceases to have
    gravitational consequences. This is itself a significant result ---
    it tells us where the Jacobson mechanism \emph{stops}, which is as
    important as knowing where it operates. Constrains a broad class of
    quantum gravity phenomenology proposals.
\end{itemize}

Either outcome represents \textbf{significant progress in fundamental
physics}. This is the hallmark of a well-designed experiment: it cannot
produce an uninterpretable result.
